\subsection{Notation}\label{sets_sec}

The notation used in the model is summarized in the Tables \ref{sets} - \ref{prices}.
Table \ref{sets} defines the various dimensions which characterize an instance of the model, including the set of sectors/commodities ($i,j$), the set of regions ($r,s)$, the set of factors of production ($f$).
Set $g$ combines the production sectors $i$ and private and public consumption demand (indices \texttt{"C"} and \texttt{"G"}) and investment demand (index \texttt{"I"}).
It allows for a much tighter formulation of the model as they can all be conceived of ``goods'' produced in similar fashion.
To simplify the exposition of the model, however, we describe private consumption, public consumption, and investment demand as stand alone components.

\renewcommand{\arraystretch}{1}
\renewcommand{\tabcolsep}{1mm}
\begin{table}[t!]
\begin{center}
\caption{Definitions of set indices.}
\label{sets}
\small
\begin{tabular}{lp{4in}}
\toprule
Set &Definition  \\
\addlinespace
\midrule
\addlinespace
$i,j$ & Sectors, an aggregation of the 65 sectors in the \GTAP 12 Data Base \\
$g$ & Production sectors $i$, plus private consumption \texttt{"C"},
public demand \texttt{"G"}, and investment \texttt{"I"} \\
$r,s$ & Regions, an aggregation of the 163 regions in the \GTAP 12 Data Base\\
$f$ & Factors of production (consisting of
\textit{mobile factors}, $f \in m_f$: five labor categories
(officials, managers, and professionals; technicians and associate
professionals; clerks; service and shop workers; agricultural and other
low-skilled workers) and capital, and \textit{sector-specific} (sluggish)
factors, $f \in s_f$: agricultural land and natural resources; in the
\texttt{g20\_macro} aggregation, the five labor categories are
aggregated to a single labor input)\\ \addlinespace \bottomrule
\end{tabular}
\end{center}
\vspace{0.4in}
\end{table}

Table \ref{activities} defines the primal variables (activity levels) which characterize an equilibrium.
The model determines values of all the variables except international capital flows, a parameter which would be determined endogenously in an intertemporal model.
Table \ref{activities} also displays the concordance between the variables and their \GAMS equivalents.

\renewcommand{\arraystretch}{1}
\renewcommand{\tabcolsep}{1mm}
\begin{table}[t!]
\begin{center}
\caption{Definitions of activity levels (quantity variables).}
\label{activities}
\small
\begin{tabular}{lllr}
\toprule
Variable &Definition  & GAMS variable & Benchmark (\GTAP) value \\
\addlinespace
\midrule
\addlinespace
$Y_{ir}$&       Production & \texttt{Y(i,r)} &\texttt{vom(i,r)}\\
\\
$C_{r}$&        Discretionary consumption & \texttt{Y("c",r)} &\texttt{vom("c",r)}\\
\\
$G_{r}$&        Aggregate public & \texttt{Y("g",r)}&\texttt{vom("g",r)} \\
\\
$I_{r}$&        Aggregate investment & \texttt{Y("i",r)}&\texttt{vom("i",r)} \\
\\
$M_{ir}$&       Aggregate imports & \texttt{M(i,r)}&\texttt{vim(i,r)}\\
\\
$X_{ir}$&       Trade flows with rest of world regions &
\texttt{X(i,r)}&\texttt{vem(i,r)}\\
\\
$\textit{FT}_{fr}$&      Factor transformation &
\texttt{FT(f,r)}&\texttt{evom(f,r)}\\
\\
$\textit{YT}_j$  &       International transport services & \texttt{YT(j)}&\texttt{vtw(j)}\\
\addlinespace
\bottomrule
\end{tabular}
\end{center}
\vspace{0.4in}
\end{table}

Table \ref{prices} defines the relative price variables for goods and factors in the model.
As is the case in any Shoven-Whalley \CGE model, the equilibrium conditions determine \textit{relative} rather than \textit{nominal} prices.

\renewcommand{\arraystretch}{1}
\renewcommand{\tabcolsep}{1mm}
\begin{table}[t!]
\begin{center}
\caption{Definitions of price variables.}
\label{prices}
\small
\begin{tabular}{lp{3in}r}
\toprule
Variable &Definition & GAMS variable\\
\addlinespace
\midrule
\addlinespace
$p^C_{r}$      &       Final demand price index for the Cobb-Douglas
demand system &\texttt{P("C",r)}\\
\\
$p^G_{r}$      &       Public provision price index  &\texttt{P("G",r)}\\
\\
$p^I_{r}$      &       Investment price index &\texttt{P("I",r)}\\
\\
$p^Y_{ir}$     &       Supply price, gross of indirect producer
taxes (when $\eta^{DX}_{i} < \infty$, the supply price is
replaced by and export price, $p^X_{ir}$ and domestic supply price,
$p^D_{ir}$) &\texttt{P(i,r)}\\
\\
$p^M_{ir}$     &       Import price, gross of export
taxes and tariffs.&\texttt{PM(i,r)}\\
\\
$p^E_{ir}$     &       Export price from endogenous or exogenous
regions.&\texttt{PE(i,r)}\\
\\
$p^T_j$         &       Marginal cost of transport services
&\texttt{PT(j)}\\
\\
$p^F_{f,r}$     &       Price of mobile primary factors ($m_f$
include labor, land, and resources) &\texttt{PF(f,r)}\\
\\
$p^S_{s_f,ir}$     &       Price of sector-specific primary
factors&\texttt{PS(sf,i,j)}\\
\addlinespace
\bottomrule
\end{tabular}
\end{center}
\vspace{0.4in}
\end{table}

While the core \GTAPINGAMS model is base on Cobb-Douglas final demand, the canonical models incorporate logic for two additional demand systems: the \LES and the \CDE demand system.
In particular, the \LES representation distinguishes between subsistence demand and discretionary demand, and we define the additional variables in Table \ref{dsysprices}.
Similarly, we define an Armington composite price index for products entering \CDE final demand.
(Details of these demand systems are provided in Appendices C and D of \citet{LanzRutherford2016}.)

\begin{table}[t!]
\begin{center}
\caption{Additional variables for \LES and \CDE demand systems.}
\label{dsysprices}
\small
\begin{tabular}{lp{3in}r}
\toprule
Variable &Definition & GAMS variable\\
\addlinespace
\midrule
\addlinespace
$SD_{r}$&       Subsistence demand & \texttt{Y("sd",r)} \\
\\
$DD_{r}$&       Discretionary demand & \texttt{Y("dd",r)} \\
\\
$p^{DD}_{r}$   &       Discretionary demand price index &\texttt{P("dd",r)}\\
\\
$p^{SD}_{r}$   &       Subsistence demand price index &\texttt{P("sd",r)}\\
\\
$p^A_{iCr}$     &      Armington composite price representing the
market price of goods entering \CDE demand.&\texttt{PA(i,r)}\\
\addlinespace
\bottomrule
\end{tabular}
\end{center}
\vspace{0.4in}
\end{table}

Finally, Table \ref{taxes} reports the definition of tax and subsidy rates applied in the model, both in terms of the notation employed to describe the model and that used in the \GAMS code.
Note that revenues from taxes and subsidy expenditures do not appear as explicit variables in the \GTAP Data Base and are defined on the basis of expenditures and tax rates.
We come back to this below.

\renewcommand{\arraystretch}{1}
\renewcommand{\tabcolsep}{1mm}
\begin{table}[t!]
\begin{center}
\caption{Tax and subsidy rates (net basis unless noted).}
\label{taxes}
\small
\begin{tabular}{lp{6cm}lr}
\toprule
Parameter & Definition &  & \GAMS Parameter \\
\addlinespace
\midrule
\addlinespace
$t^o_{ir}$ & Output taxes (gross basis)& & \texttt{rto(i,r)}    \\
\\
$t^{f}_{fir}$   &Factor taxes &  & \texttt{rtf(f,i,r)}   \\
\\
$t^{fd}_{ijr}$ & Intermediate input taxes &Domestic&\texttt{rtfd(i,j,r)}      \\
$t^{fi}_{ijr}$ &                          &Imported& \texttt{rtfi(i,j,r)}      \\
\\
$t^{pd}_{ir}$ & Consumption taxes  & Domestic& \texttt{rtfd(i,"C",r)}     \\
$t^{pi}_{ir}$ &                 &Imported& \texttt{rtfi(i,"C",r)} \\
\\
$t^{gd}_{ir}$ & Public demand taxes &Domestic & \texttt{rtfd(i,"G",r)}   \\
$t^{gi}_{ir}$&                 &Imported & \texttt{rtfi(i,"G",r)} \\
\\
$t^{gd}_{ir}$ & Investment demand taxes &Domestic & \texttt{rtfd(i,"I",r)}   \\
$t^{gi}_{ir}$&                 &Imported& \texttt{rtfi(i,"I",r)} \\
\\
$t^{xs}_{isr}$  &Export subsidies&  & \texttt{rtxs(i,s,r)}   \\
\\
$t^{ms}_{isr}$  &Import tariffs& & \texttt{rtms(i,s,r)}\\
\addlinespace
\bottomrule
\end{tabular}
\end{center}
\vspace{0.4in}
\end{table}

\subsection[Accounting identities]{Benchmark data structure and accounting identities}

The economic structure underlying the \GTAP dataset and model is illustrated in Figure \ref{fig:region}.
Symbols in this flow chart correspond to variables in the economic model (see Table \ref{activities}): $Y_{ir}$ is the production of good $i$ in region $r$, $C_r$, $I_r$, and $G_r$ portray private consumption, investment, and public demand, respectively, $M_{ir}$ portrays the import of good $i$ into region $r$, $\textit{RA}_r$ stand for representative consumers, and $FT_{s_{f}r}$ is the activity through which the set of sector-specific factors of production ($s_f$) are allocated to individual sectors.
Further, solid lines represent commodity and factor market flows, while dotted lines indicate tax revenues and transfers.

\begin{sidewaysfigure}
\begin{center}
\includegraphics{jgea_new/figures/gtapflowchart.tikz}
\caption{Regional economic structure.}
\label{fig:region}
\end{center}
\end{sidewaysfigure}

Domestic and imported goods markets are represented by horizontal lines at the top of the figure.
Domestic production ($\textit{vom}_{ir}$) is distributed to exports ($\textit{vxmd}_{irs}$), international transportation services ($\textit{vst}_{ir}$), intermediate demand ($\textit{vdfm}_{ijr}$), household consumption ($\textit{vdfm}_{iCr}$), investment ($\textit{vdfm}_{iIr}$), and government consumption ($\textit{vdfm}_{iGr}$).\footnote{Recall that in the \GAMS implementation of the model the index \texttt{"g"} includes all sectors represented in the model plus private consumption \texttt{"C"}, public demand \texttt{"G"}, and investment \texttt{"I"}.
(See Table \ref{sets}.)
For the LES demand representation, it also includes price indices for discretionary demand \texttt{"dd"} and subsistence demand \texttt{"sd"}. } The accounting identity in the \GTAP dataset is thus:
$$\textit{vom}_{ir} = \sum_s \textit{vxmd}_{irs} + \textit{vst}_{ir}
+ \sum_j \textit{vdfm}_{ijr} + \textit{vdfm}_{iCr} + \textit{vdfm}_{iIr} + \textit{vdfm}_{iGr}\,,$$ where $j$ indexes all goods.
Similarly, imported goods (with aggregate value $\textit{vim}_{ir}$) enter intermediate demand ($\textit{vifm}_{ijr}$), private consumption ($\textit{vifm}_{iCr} + \textit{vifm}_{i\textit{SD}r}$), and public consumption ($\textit{vifm}_{iGr}$).
The accounting identity for these flows is thus:
$$\textit{vim}_{ir} = \sum_j \textit{vifm}_{ijr} +
\textit{vifm}_{iCr} + \textit{vifm}_{i\textit{SD}r} + \textit{vifm}_{iGr}\,.$$

Inputs to production of good $i$ ($Y_{ir}$) include intermediate inputs (domestic $\textit{vdfm}_{ijr}$ and imported $\textit{vifm}_{ijr}$), mobile factors of production ($\textit{vfm}_{m_{f},ir}$, where $m_{f}$ is a subset of the set $f$ designating all factors of production), and sector-specific factors of production ($\textit{vfm}_{s_{f},ir}, s_f \subset f$).
Factor market equilibrium is given by an identity relating the value of factor payments to factor income:
$$\sum_i \textit{vfm}_{fir} = \textit{evom}_{fr}\,,$$ and factor
earnings accrue to households.

International market clearance conditions require that region $r$ exports of good $i$ ($\textit{vxmd}_{ir}$ at the top of the figure) equal the imports of the same good from the same region summed across all trading partners ($\textit{vxmd}_{isr}$ at the bottom of the figure): $$\textit{vxm}_{ir} = \sum_s \textit{vxmd}_{irs}\,,$$ where $s$, an alias for $r$, indexes regions.
Likewise, market clearance conditions apply for international transportation services.
The supply-demand balance in the market for transportation service $j$ requires that the sum across all regions of service exports ($\textit{vst}_{ir}$, at the top of the figure) equals the sum across all bilateral trade flows of service inputs ($\textit{vtwr}_{jisr}$ at the bottom of the figure):
$$\sum_r \textit{vst}_{jr} = \sum_{isr} \textit{vtwr}_{jisr}$$

Turning to tax revenues and transfers, shown as dotted lines in figure \ref{fig:region}, flows labeled with $\mathcal{R}$ correspond to tax revenues.
For each country, tax flows consist of indirect taxes on production/exports of each good ($\mathcal{R}^Y_{ir}$), on consumption ($\mathcal{R}^C_r$), on public demand ($\mathcal{R}^G_r$), and on imports ($\mathcal{R}^M_{ir}$).
The regional budget constraint thus relates tax payments ($\mathcal{R}^Y_{ir}, \mathcal{R}^C_r, \mathcal{R}^G_r, \mathcal{R}^M_{ir}$), factor income ($\textit{evom}_{fr}$), and the current account deficit (i.e., net transfers from abroad, $\textit{vb}_r$) to total private consumption expenditure $\textit{vom}_{Cr}$, total public consumption expenditure $\textit{vom}_{Gr}$, and total investment $\textit{vom}_{Ir}$, yielding:
$$\textit{vom}_{Cr} + \textit{vom}_{Gr} + \textit{vom}_{Ir} = \sum_f
\textit{evom}_{fr} + \sum_i \mathcal{R}^Y_{ir} + \mathcal{R}^C_r + \mathcal{R}^G_r + \sum_i \mathcal{R}^M_{ir} + \textit{vb}_r\,.
$$

To this point we have outlined two types of consistency conditions which are part of any \GTAP Data Base: market clearance (supply = demand for all goods and factors), and income balance (net income = net expenditure).
A third set of identities involve net operating profits by all sectors in the economy.
In the core \GTAP model ``production'' takes place under conditions of perfect competition with constant returns to scale, hence there are no excess profits, and the cost of inputs must equal the value of outputs.
This condition applies for each production sector:

\begin{itemize}
\item[$Y_{ir}$:] $\sum_f \textit{vfm}_{fir} + \sum_j \left(\textit{vdfm}_{ijr} + \textit{vifm}_{ijr} \right) + \mathcal{R}^Y_{ir} = \textit{vom}_{ir}\,,$
\item[$C_{r}$:] $\sum_i \left(\textit{vdfm}_{iCr} + \textit{vifm}_{iCr} \right) + \mathcal{R}^C_{ir} = \textit{vom}_{Cr}\,,$
\item[$I_{r}$:] $\sum_i \textit{vdfm}_{iIr} = \textit{vom}_{Ir}$
\item[$G_{r}$:] $\sum_i \left(\textit{vdfm}_{iGr} + \textit{vifm}_{iGr} \right) + \mathcal{R}^G_{ir} = \textit{vom}_{Gr}\,,$
\item[$M_{ir}$:] $\sum_s \left(\textit{vxmd}_{isr} + \sum_j \textit{vtwr}_{jisr} \right) + \mathcal{R}^M_{ir} = \textit{vim}_{ir}$
\item[$FT_{fr}$:] $\textit{evom}_{fr} = \sum_i \textit{vfm}_{fir}$
\end{itemize}

\subsection[Optimization problems]{Decentralized optimization problems}

The benchmark identities presented in the previous section indicate the market clearance, zero profit, and income balance conditions which define the \GTAP model.
The displayed equations do not, however, characterize the behavior of agents in the model.
In a competitive equilibrium setting, the standard assumption of optimizing atomistic agents applies for both producers and consumers.
This section lays out the optimization problem of each component in the model, and thereby provides the structure of production technology (production functions) and preferences (characterizing final demand), as well as the representation of trade.
The treatment of the rest of the world, which distinguishes the \GMR model from the small open economy closures, is described in Appendix \ref{app:closures}.

Note that in order to simplify notation, we denote decision variables corresponding to the benchmark data structures with the initial ``v'' replaced by ``d.''
Hence, while $vdfm_{jir}$ represents benchmark data on intermediate demand for good $j$ in the production of good $i$ in region $r$, $ddfm_{jir}$ represents the corresponding decision variable in the equilibrium model.
This approach to the scaling of variables is consistent with the \GAMS code, and it provides a flexible and transparent approach with respect to the calibration of activity variables.

\subsubsection{Production technology}

Starting with producers, profit maximization in the constant returns to scale setting is equivalent to cost minimization subject to technical constraints.
For sector $Y_{ir}$ we characterize input choices as though they arose from minimization of unit costs:
\begin{eqnarray}\label{profitmax}
\min_{\textit{ddfm,dfm,difm}} \quad & c^D_{ir} + c^M_{ir} +
c^F_{ir} & \\
&s.t.&c^D_{ir} = \sum_j p^Y_{jr} (1 + t^{fd}_{jir})
\textit{ddfm}_{jir} \nonumber\\
&&c^M_{ir} = \sum_j p^M_{jr} (1 + t^{fi}_{jir}) \textit{difm}_{jir}\nonumber\\
&&c^F_{ir} = \sum_f (p^F_{fr}|_{f \in m_f} + p^S_{fir}|_{f \in s_f}) (1 +
t^f_{fir} ) \textit{dfm}_{fir}\nonumber\\
&&F_{ir}(\textit{ddfm,difm,dfm}) = Y_{ir}\nonumber
\end{eqnarray}
where $F(\cdot)$ represents the production function, which is described by a nested constant-elasticity-of-substitution (CES) form, with structure displayed in Figure \ref{prodfunY}.

\begin{figure}[t!]
\centering
\includegraphics{jgea_new/figures/prodfunY.tikz}
\vspace{0.1in}
\caption{CES nesting structure for production function $Y_{ir}=F_{ir}(\textit{ddfm,difm,dfm})$.}
\label{prodfunY}
\vspace{0.3in}
\end{figure}

In the figure, $\sigma$ values in different nests represent substitution elasticities between inputs, with $\sigma^D_i = \textit{esubd}_i$ measuring substitution possibility between intermediate inputs produced domestically and imported from abroad which are similarly a composite of imports from varieties from different regions.
$\textit{esubva}_i$ represents the elasticity of substitution between primary inputs in the value added nest.
(Both of these parameters are provided in the \GTAP Data Base \citep[see][]{HertelvanderMensbrugghe2015}.)
Note further that the specific source of tax revenue is indicated in this figure, consisting of output taxes, taxes on intermediate inputs, and taxes on factor demands, all of which are applied on an ad-valorem basis.

One important modelling choice concerns the differentiation of output for domestic and export markets.
In particular, the supply of goods to domestic and export markets are portrayed as arising from the following profit-maximization problem:
\begin{eqnarray}\label{cet_dx}
\max_{\textit{dxm},\textit{ddm}} \quad && p^D_{ir} \textit{ddm}_{ir}
+ p^X_{ir} \textit{dxm}_{ir}
\\
&s.t.& G_{ir}(\textit{ddm}_{ir},\textit{dxm}_{ir}) = Y_{ir}
\end{eqnarray}
where $G$ is the \CET function with structure illustrated in Figure \ref{prodfunDX}.
In the standard \GTAP model, transformation elasticities $\eta^{DX}_{ir} = \textit{etrndx}_{i}$ are infinite.
This implies that the supply price of output $p^Y_{jr}$ is the same if output is used locally or in a different region.
With $\textit{etrndx}_{i} < \infty$, the output price may differ when supplied domestically or abroad.
In this case, $p^Y_{jr}$ in equation \eqref{profitmax} is replaced by $p^D_{jr}$.
In \GTAPINGAMS, the choice of $\textit{etrndx}_{i}$ is governed by the parameter \texttt{thetadx} (see Appendix \ref{app:thetadx}) and applies to both the \GMR and the small open economy closures: \texttt{thetadx}$=1$ yields $\textit{etrndx}_{i}=\infty$ (the \texttt{gtap} model version of this paper), \texttt{thetadx}$=0.5$ yields a finite elasticity (the \texttt{gtapingams} model version, the default of \citealp{LanzRutherford2016}).

\begin{figure}[t!]
\centering
\includegraphics{jgea_new/figures/prodfunDX.tikz}
\vspace{0.1in}
\caption{CET transformation between domestic and export markets $\textit{Y}_{ir} = G_{ir}(\textit{dxm},\textit{ddm})$.}
\label{prodfunDX}
\vspace{0.3in}
\end{figure}

\subsubsection{Preferences and final demand}

Private consumption consistent with utility maximization is portrayed by minimization of the cost of a given level of aggregate consumption:
\begin{eqnarray}\label{utils}
\min_{\textit{ddfm}_{iCr},\textit{difm}_{iCr}} \quad &&\sum_i p^Y_{ir}
(1 + t^{pd}_{ir}) \textit{ddfm}_{iCr} + p^M_{ir} (1 +t^{pi}_{ir})
\textit{difm}_{iCr} \\
&s.t.& H_r(\textit{ddfm}_{iCr},\textit{difm}_{iCr}) = C_{ir} \nonumber
\end{eqnarray}
where $H_r$ represents final demand from the representative consumer.

Final demand in the core model is characterized by Cobb-Douglas preferences.
Alternative specifications included in the model logic include both the \LES and \CDE expenditure.\footnote{The \LES model involves subsistence and discretionary expenditure, which are represented through composite price indices representing \CES aggregates of both domestic and imported inputs, \texttt{P("DD",r)} and \texttt{P("SD",r)}, respectively.
The \CDE model requires the introduction of trade aggregation activities (\texttt{A(i,r)}) and associated Armington composite price indices, \texttt{PA(i,r)}.}
The nested discretionary and subsistence demand functions are is displayed in Figure \ref{prodfunC}.

\begin{figure}[t!]
\centering
\includegraphics{jgea_new/figures/prodfunC.tikz}
\vspace{0.1in}
\caption{LES nesting for discretionary and subsistence consumption.}
\label{prodfunC}
\vspace{0.3in}
\end{figure}

\subsubsection{Government and public consumption}

Public consumption in the model is represented as a fixed coefficient (Leontief) aggregation of domestic-import composites.
This formulation introduces substitution at the second level between domestic and imported inputs while holding sectoral commodity aggregates constant.
Figure \ref{prodfunG} illustrates the functional form.

\begin{figure}[t!]
\centering
\includegraphics{jgea_new/figures/prodfunG.tikz}
\vspace{0.1in}
\caption{Nesting structure for public consumption $G_r = G_r(\textit{ddfm}_{iGr},\textit{difm}_{iGr})$.}
\label{prodfunG}
\vspace{0.3in}
\end{figure}

\subsubsection{International trade}

The choice among imports from different trading partners is based on Armington's idea of regionally differentiated products.
The following cost minimization problem formalizes this choice:
\begin{eqnarray}\label{optA}
\min_{\textit{dxmd,dtwr}} \quad &&\sum_s (1+t^{ms}_{isr}) \left(p^Y_{is}
(1-t^{xs}_{isr}) \textit{dxmd}_{isr} +
\sum_j p^T_j \textit{dtwr}_{jisr} \right)\\
&s.t.&A_{ir}(\textit{dxmd,dtwr}) = M_{ir} \nonumber
\end{eqnarray}
where $A$ is the import aggregation function, described by the nested \CES-Leontief function shown in Figure \ref{prodfunA}.
When $\textit{etrndx}_{i}<\infty$, $p^Y_{is}$ in (\ref{optA}) is replaced by $p^X_{is}$.

\begin{figure}[t!]
\centering
\includegraphics{jgea_new/figures/prodfunA.tikz}
\vspace{0.1in}
\caption{Armington aggregation of traded goods $M_{ir} = A_{ir}(\textit{dxmd,dtwr})$.}
\label{prodfunA}
\vspace{0.3in}
\end{figure}

Note that transportation services enter on a proportional basis with imports from different countries, reflecting differences in unit transportation margins across different goods and trading partners.
Therefore, substitution at the top level in an Armington composite involves trading off of both imported goods and associated transportation services.
Trade flows are subject to export subsidies and import tariffs, with subsidies paid by government in the exporting region, and tariffs collected by government in the importing region.

The provision of international transportation services is modeled through an aggregation of transportation services exported from countries throughout the world.
More specifically, we consider the following cost minimization problem for the aggregation of transportation services:
$$\min_{\textit{dst}} \quad \sum_r p^Y_{ir} \textit{dst}_{ir} \quad
\text{s.t.}
\quad T_i(\textit{dst}) = \textit{YT}_i$$ where the aggregation function $T_i$ combines transport service exports from multiple regions.
The functional form which aggregates services from different regions is illustrated in Figure \ref{prodfunYT}.

\begin{figure}[t!]
\centering
\includegraphics{jgea_new/figures/prodfunYT.tikz}
\vspace{0.1in}
\caption{International transportation services aggregator $\textit{YT}_j = T_j(\textit{dst})$.}
\label{prodfunYT}
\vspace{0.3in}
\end{figure}

\subsubsection{Supply of sector-specific factors}

Land and natural resources are portrayed as sector-specific factors of production.
These are supplied through constant-elasticity-of-transformation (CET) production function which allocates composite factors to sectoral markets.
Formally, the supply of sectoral factors of production is modeled through the following profit-maximization problem:
\begin{eqnarray}\label{cet_input}
\max_{\textit{dfm}} \quad && \sum_j \textit{dfm}_{s_{f},jr} p^S_{s_{f},jr}\\
&s.t.& \Gamma_{s_{f},r}(\textit{dfm}) = \textit{evom}_{s_{f},r}
\end{eqnarray}
where $\Gamma$ is the \CET function with structure illustrated in Figure \ref{prodfunFT}.
Note that in the figure $\eta$ represents transformation elasticities provided in the \GTAP Data Base.

\begin{figure}[t!]
\centering
\includegraphics{jgea_new/figures/prodfunFT.tikz}
\vspace{0.1in}
\caption{Sector-specific factor \CET transformation function $\textit{FT}_{s_{f}} = \Gamma_{s_{f}}(\textit{dfm})$.}
\label{prodfunFT}
\vspace{0.3in}
\end{figure}

\subsection[Equilibrium conditions]{Equilibrium Conditions}\label{equil}

An Arrow-Debreu model concerns itself with the interactions of decentralized decisions by consumers and producers in markets.
\citet{Lars} proposed a representation of this class of models in which two types of equations define an equilibrium: zero profit and market clearance.
The corresponding variables defining an equilibrium are activity levels (for constant-returns-to-scale firms) and commodity prices.\footnote{Under a maintained assumption of perfect competition, \citet{Lars} may characterize technology as constant-returns-to-scale without loss of generality.
Specifically, decreasing returns are accommodated through introduction of a specific factor, while increasing returns are inconsistent with the assumption of perfect competition.
Note that in this environment zero excess profit is consistent with free entry for atomistic firms producing an identical product.}
Here we extend Mathiesen's framework with a third class of variables corresponding to consumer income levels.
Commodity markets encompass primary endowments of households with producer outputs.
In equilibrium the aggregate supply of each good must be at least as great as total intermediate and final demand.
Initial endowments are exogenous.
Producer supplies and demands are defined by producer activity levels and relative prices.
Final demands are determined by market prices.

Economists who have worked with conventional textbook equilibrium models can find Mathiesen's framework to be somewhat opaque because many quantity variables need not be explicitly specified in the model.
Variables such as final demand by consumers, factor demands by producers, and commodity supplies by producers, are defined implicitly in Mathiesen's model.
For example, given equilibrium prices for primary factors, consumer incomes can be computed, and given income and goods prices, consumers' demands can then be determined.
The consumer demand functions are written down in order to define an equilibrium, but quantities demanded need not appear in the model as separate variables.
The same is true of inputs or outputs from the production process: relative prices determine conditional demand, and conditional demand times the activity level represents market demand.
Omitting decisions variables and suppressing definitional equations corresponding to intermediate and final demand provides significant computational advantages at the cost of a somewhat more complicated model statement.

In the following, we detail (i) zero profit conditions, (ii) market clearance conditions, and (iii) income balance conditions, which in the present case is equivalent to the regional budget constraint.
These three sets of conditions form the basic system of equation to be solved.
Note that the actual code for the model is implemented both in the algebraic mixed complementarity format \citep[\GAMS/\MCP, see][]{Rutherford1995} and through the more compact formulation afforded by the Mathematical Programming System for General Equilibruim (\MPSGE) syntax \citep{Rutherford1999}.

\subsubsection[Zero profit]{Zero profit (arbitrage) conditions}

All production activities in the model are represented by constant-returns-to-scale technologies, and markets are assumed to operate competitively with free entry and exit.
As a consequence, equilibrium profits are driven to zero and the price of output reflects the cost of inputs.
The following sets of equations relating output price to marginal cost are part of the definition of an equilibrium.\footnote{To retain consistency with the \MCP format, we express zero profit conditions as ``oriented equations,'' with marginal (=average) cost on the \LHS and marginal (=average) revenue on the \RHS.}

The calculation of unit cost and unit revenue functions involves the definition of a number of ancillary variables (that do not appear in the \GAMS code as explicit choice variables).
In the following we define ancillary variables in un-numbered equations, indicating that these variables are ``optional'' in the sense that they may be substituted out of the non-linear system of equations.
Moreover, we use the symbol $\theta$ to portray value shares from the base year data.
In most cases subscripts on these value shares are omitted in order to economize on notation.
Finally, to denote benchmark values we use an overline, so that $\overline{t}^{pd}_{ir}$ represents the benchmark value of ${t}^{pd}_{ir}$.

\paragraph{\bf Sectoral production (\texttt{Y(j,r)})}

Sectoral production combines intermediate inputs with a value-added nest combining primary inputs (see Figure \ref{prodfunY}).
The unit cost of value-added is a \CES composite of skilled and unskilled labor, land, resources, and capital inputs to production, gross of taxes.
Factor inputs may be sector-specific or mobile across sectors:
$$p^{pf}_{fjr} = \left\{\begin{array}{lr}
p^F_{fr} \frac{(1+t^f_{fjr})}{1+\overline{t}^f_{fjr}} & f \in m_f\\
p^S_{fjr} \frac{(1+t^f_{fjr})}{1+\overline{t}^f_{fjr}} & f \in s_f
\end{array} \right.$$
and the unit cost function is given by:
$$c^f_{jr} = \left( \sum_f \theta_f \left( p^{pf}_{fjr}
\right)^{1-\sigma} \right)^{1/(1-\sigma)},$$ in which $\theta_f$ represents the factor share of value added [{\small \texttt{theta\_vfm(f,g,r)}}].

The user cost of intermediate inputs differs from the market price due to the presence of taxes on intermediate inputs:
$$p^d_{ijr} = p^Y_{ir} \frac{1+t^{fd}_{ijr}}{1+\overline{t}^{fd}_{ijr}}$$
$$p^i_{ijr} = p^M_{ir}
\frac{1+t^{fi}_{ijr}}{1+\overline{t}^{fi}_{ijr}}$$ A \CES cost function describes the minimum cost of a bundle of domestic and imported inputs to production, based on benchmark value shares and an elasticity of substitution $\sigma [\texttt{esubd(i)}]$:
$$c^i_{ijr} = \left( \theta (p^d_{ijr})^{1-\sigma} +
             (1-\theta) (p^i_{ijr})^{1-\sigma} \right)^{1/(1-\sigma)},$$ in which $\theta$ represents the domestic share of the Armington composite \hfill \break [{\small \texttt{theta\_vdfm(i,j,r)}}].

Unit cost of sectoral output is then a Leontief (linear) composite of the costs of intermediate and value-added composite inputs, based on base-year value shares:
$$c^Y_{jr} = \sum_i \theta_i c^i_{ijr} + \theta^f c^f_{jr},$$
in which $\theta_i$ represents the cost share of intermediate input $i$ [{\small \texttt{theta\_cm0(i,j,r)}}]; and $\theta^f$ represents the cost share of value added in sectoral output [{\small \texttt{theta\_cf0(j,r)}}].

Having formulated the unit cost function, it is possible to compactly portray the zero profit condition for $y_{jr}$.
In equilibrium, the marginal cost of supply equals the market price, net of taxes:
\begin{equation}
c^Y_{jr} = p^Y_{jr} \frac{1-t^o_{jr}}{1-\overline{t}^o_{jr}}
\end{equation}
When $\textit{etrndx}_{i}<\infty$, production for domestic and export markets is differentiated, and we replace $p^Y_{jr}$ by a unit revenue function,
$$r^Y_{jr} = \left( \theta (p^D_{jr})^{1+\eta} +
(1-\theta) (p^X_{jr})^{1+\eta} \right)^{1/(1+\eta)},$$ in which $\theta$ represents the export share of output [{\small \texttt{theta\_vxm(j,r)}}].

\paragraph{\bf Consumer demand (\texttt{Y("c",r)})}

In previous versions of \GTAPINGAMS the consumer price index represented a Cobb-Douglas demand system with an Armington substitution between domestic and imported goods prices gross of tax.
This model is retained as one option in the new model, price indices for domestic and imported goods are given by:
$$p^{dc}_{ir} = p^Y_{ir} \frac{1+t^{pd}_{ir}}{1+\overline{t}^{pd}_{ir}},$$
and
$$p^{ic}_{ir} = p^M_{ir} \frac{1+t^{pi}_{ir}}{1+\overline{t}^{pi}_{ir}},$$
and the Armington composite price of good $i$ is a \CES composite price defined over these indices:
$$p^c_{ir} = \left(\theta (p^{dc}_{ir})^{1-\sigma} +
                (1-\theta) (p^{ic}_{ir})^{1-\sigma} \right)^{1/(1-\sigma)},$$ in which $\theta$ represents the domestic demand share [{\small \texttt{theta\_vdfm(i,"c",r)}}].

In addition to the Cobb-Douglas representation of the final demand system, our core model includes a \LES, which combines a Cobb-Douglas discretionary demand with a Leontief subsistence demand activity.
The added complexity of the \LES specification permits the model to be calibrated both to the income and \textit{average} price elasticities of demand.
(See Appendix C of \citet{LanzRutherford2016} for details.)
The price indices for discretionary and subsistence demand are defined by:
\begin{equation}
p^{DD}_r = \prod_i (p^c_{ir})^{\theta_i},
\end{equation}
in which $\theta_i$ represents the discretionary value share [{\small \texttt{theta\_cm0(i,"dd",r)}}].

The price index for subsistence demand is defined by a Leontief cost function:
\begin{equation}
p^{SD}_r = \sum_i \theta_i p^c_{ir},
\end{equation}
in which $\theta_i$ represents the commodity $i$ share of subsistence demand \hfill \break [{\small \texttt{theta\_cm0(i,"sd",r)}}].

Finally, \GAMS code includes a third representation of final demand in the form of the \CDE demand system \citep{Hanoch1975}.
This representation, details of which are presented in Appendix D of \citet{LanzRutherford2016}, has been part of the \GTAP model in \textsc{gempack} for many years, and our implementation follows those of \citet{HertelvanderMensbrugghe2015} and \citet{Chen}.
Similar to the LES demand system, this approach permits calibrating empirical evidence on both own-price and income elasticities.

\paragraph{\bf Government demand (\texttt{Y("g",r)})}

Public expenditure is a fixed-coefficient aggregate of Armington composite goods.
Within each composite domestic and imported goods trade off with a constant elasticity of substitution.
The unit price indices for domestic and imported goods are given by:
$$p^{dg}_{ir} = p^Y_{ir} \frac{1+t^{dg}_{ir}}{1+\overline{t}^{gd}_{ir}} $$
and
$$p^{ig}_{ir} = p^M_{ir} \frac{1+t^{ig}_{ir}}{1+\overline{t}^{ig}_{ir}}$$
The composite price of the $i$th good is then:
$$pg_{ir} = \left( \theta (p^{dg}_{ir})^{1-\sigma} + (1-\theta)
(p^{ig}_{ir})^{1-\sigma} \right)^{1/(1-\sigma)},$$ in which $\theta$ represents the domestic share of public demand [{\small \texttt{theta\_vdfm(i,"g",r)}}].

The unit cost of public services ($G_{r}$) is defined by the Leontief cost coefficients:
\begin{equation}
\sum_i \theta_i pg_{ir}  = p^G_r,
\end{equation}
in which $\theta_i$ represents the value share of commodity $i$ in public expenditure \hfill \break [{\small \texttt{theta\_cm0(i,"g",r)}}].

\paragraph{\bf Aggregate imports (\texttt{M(i,r)})}

An import cost index applies export taxes, trade and transport margins, and import tariffs to the producer supply prices in exporting regions:
$$py^m_{isr} = p^Y_{is}
\frac{(1-t^{xs}_{isr})(1+t^{ms}_{isr})} {(1-\overline{t}^{xs}_{isr})(1+\overline{t}^{ms}_{isr})}\,.$$ When $\textit{etrndx}_{i}<\infty$, the supply price $p^Y_{is}$ is replaced by export price $p^X_{is}$.
The unit price of transportation services is given by: $$pt^m_{jisr} = p^T_j \frac{1+t^{ms}_{isr}}{1+\overline{t}^{ms}_{isr}}\,.$$

Transportation margins enter as fixed coefficients with bilateral trade flows, so the unit delivered price is a convex combination of the unit prices with weights corresponding to base year value shares:
$$pyt^m_{isr} = \theta py^m_{isr} + \sum_j \theta^T_j pt^m_{jisr},$$
in which $\theta$ represents good $i$ share of imports [{\small \texttt{theta\_vxmd(i,s,r)}}]; and $\theta^T_j$ represents the value share of transportation service $j$ in the import price \hfill \break [{\small \texttt{theta\_vtwr(j,i,s,r)}}].

Having formed a price index for bilateral imports from region $s$ to region $r$, the \CES cost index can be defined on the basis of value shares and the elasticity of substitution across imports from different regions, $\sigma = \textit{esubm}_i$:
$$\textit{cim}_{ir} = \left(\sum_s \theta_s (pyt^m_{isr})^{1-\sigma)}
\right)^{1/(1-\sigma)} ,$$ in which $\theta_s$ is the value share of bilateral imports from region $s$ [{\small \texttt{theta\_m(i,s,r)}}].

The import activity ($m_{ir}$) has a zero profit condition which relates the unit cost of imports to the market price of the import aggregate:
\begin{equation}
cim_{ir} = p^M_{ir}\,.
\end{equation}

\paragraph{\bf International transportation services (\texttt{YT(j)})}

For simplicity, the unit cost of a transportation service depends on the benchmark value shares of region-specific services through a Cobb-Douglas cost function.
Under perfect competition with free entry, the unit cost of international transport services equals the equilibrium market price:
\begin{equation}
\prod_r (p^Y_{jr})^{\theta_r} = p^T_j
\end{equation}
in which $\theta_r$ represents the region $r$ share of transportation service $j$ [{\small \texttt{theta\_vst(j,r)}}].
When $\textit{etrndx}_{i}<\infty$, the supply price $p^Y_{is}$ in this equation is replaced by export price $p^X_{is}$.

\paragraph{\bf Sector-specific factor transformation (\texttt{FT(f,r)})}

The unit value of sector-specific factors is defined as a \CET aggregate of returns to factor $f$ across sectors $j$:
$$\textit{pvfm}_{fr} = \left( \sum_j \theta_j ps_{fjr}^{1+\eta}
\right)^{1/(1+\eta)} \quad f \in s_{f}\,.$$ in which $\theta_j$ is the sector $j$ share of earnings for factor $f$ [{\small \texttt{theta\_evom(f,j,r)}}].

The constant elasticity of transformation frontier defines the profit-maximizing allocation of factors to individual sectors.
In equilibrium, the unit value of the aggregate factor is equal to the maximum unit earnings:
\begin{equation}
p^S_{fr} = \textit{pvfm}_{fr} \quad f \in s_f
\end{equation}

\subsubsection{Market clearance}

Supply-demand conditions apply to all goods and factors.
Benchmark demand and supply quantities appear as scale factors in many of these equations, typically multiplied by activity levels which are equal to unity in the reference equilibrium.\footnote{While not crucial for representation of the model as a nonlinear system of equations, we follow the \MCP convention in writing out the market clearance conditions.
The equations are ``oriented'', with supply variables on the \LHS and demands on the \RHS.
Hence, the sense of the equation is $supply \geq demand$.
In the core model equilibrium prices should always be positive, but in extensions of the standard model it might be quite common to introduce inequalities and complementary slackness, in which case the proper orientation of the equations is essential.
Hence, in equilibrium should the price of a good be zero, economic equilibrium is then consistent with a market in which $supply > demand$.}

\paragraph{\bf Firm output (\texttt{P(i,r)})}

Aggregate output of good $i$ in region $r$ in the reference equilibrium is \texttt{vom(i,r)}:
\begin{equation}
\label{Youtput}
Y_{ir} \textit{vom}_{ir} =
\sum_j \textit{ddfm}_{ijr} + \textit{ddfm}_{iCr} + \textit{ddfm}_{iIr}
+ \textit{ddfm}_{iGr} + \sum_s \textit{dxmd}_{irs} + \textit{dst}_{ir}
\end{equation}
where the compensated demand functions can be obtained by differentiating the unit cost functions:
\begin{eqnarray}
\textit{ddfm}_{ijr} &=& Y_{jr} \textit{vdfm}_{ijr} \left( \frac{\textit{ci}_{ijr}}{p^d_{ijr}} \right)^\sigma \nonumber\\
\textit{ddfm}_{iCr}  &=& C_r \left( \textit{vdfm}_{iCr}\frac{p^C_r}{pc_{ir}} \right) \left(\frac{p^C_{ir}}{p^{\textit{dc}}_{ir}} \right)^\sigma \nonumber\\
\textit{ddfm}_{iIr}  &=& I_r \textit{vdfm}_{iIr} \nonumber\\
\textit{ddfm}_{iGr} &=& G_r \textit{vdfm}_{iGr}\left(\frac{pg_{ir}}{p^{\textit{dg}}_{ir}} \right)^\sigma   \nonumber\\
\textit{dxmd}_{isr} &=& M_{ir} \textit{vxmd}_{isr} \left(\frac{p^M_{ir}}{\textit{pyt}^m_{isr}} \right)^\sigma \nonumber\\
\textit{dst}_{jr} &=& \textit{YT}_j  \textit{vst}_{jr} \frac{p^T_j}{p^Y_{jr}}\nonumber
\end{eqnarray}

When $\textit{etrndx}_{i}<\infty$, equation (\ref{Youtput}) is replaced by two equations, one representing the market for domestic output and another representing the market for exports:
\begin{eqnarray}
Y_{ir} \textit{ddm}_{ir} &=&
\sum_j \textit{ddfm}_{ijr} + \textit{ddfm}_{iCr} + \textit{ddfm}_{iIr}
+ \textit{ddfm}_{iGr} \label{Doutput}\\
Y_{ir} \textit{dxm}_{ir} &=&
\sum_s \textit{dxmd}_{irs} + \textit{dst}_{ir}
\,.\label{Xoutput}
\end{eqnarray}
In these expressions, the compensated domestic and export supply coefficients are given by:
\begin{eqnarray}
\textit{ddm}_{ir} &=& \textit{vdm}_{ir} \left( \frac{ r^Y_{ir} } {
p^D_{ir}} \right)^{\eta} \nonumber \\
\textit{dxm}_{ir} &=& \textit{vxm}_{ir} \left( \frac{ r^Y_{ir} } {
p^X_{ir}} \right)^{\eta}. \nonumber
\end{eqnarray}

\paragraph{\bf Private consumption (\texttt{P("c",r)})}

Consumer demand in region $r$ in the reference equilibrium is $\textit{vom}_{Cr}$ hence:
\begin{equation}
C_r \textit{vom}_{Cr}  = \frac{\textit{RA}_r}{p^C_r}
\end{equation}

\paragraph{\bf Composite imports (\texttt{PM(i,r)})}

The aggregate value of imports of good $i$ in region $r$ in the reference equilibrium is $\textit{vim}_{ir}$:
\begin{equation}
M_{ir} \textit{vim}_{ir} = \sum_j \textit{difm}_{ijr} +
\textit{difm}_{iCr} + \textit{difm}_{iGr}
\end{equation}
where compensated demand functions are given by:
\begin{eqnarray}
\textit{difm}_{ijr} &=& Y_{jr} \textit{vifm}_{ijr} \left( \frac{ci_{ijr}}{p^i_{ijr}} \right)^\sigma\nonumber\\
\textit{difm}_{iCr} &=& C_r \textit{vifm}_{iCr}  \left(\frac{pc_{ir}}{p^{ic}_{ir}}\right)^\sigma  \frac{p^C_r}{\textit{pc}_{ir}} \nonumber\\
\textit{difm}_{iGr} &=& G_r \textit{vifm}_{iGr} \left(\frac{pg_{ir}}{p^{ig}_{ir}}\right)^\sigma \nonumber
\end{eqnarray}

\paragraph{\bf Transport services (\texttt{PT(j)})}

The aggregate demand (and supply) for transport service $j$ in the benchmark equilibrium is $\textit{vtw}_j$:
\begin{equation}
\textit{YT}_j  \textit{vtw}_j = \sum_{isr} \textit{dtwr}_{jisr}
\end{equation}
where
$$\textit{dtwr}_{jisr} = M_{ir} \textit{vtwr}_{jisr}  \left(
\frac{p^M_{ir}}{\textit{pyt}^m_{isr}} \right)^\sigma\,.$$

\paragraph{\bf Primary factors (\texttt{PF(f,r)})}

The aggregate demand (and supply) of primary factor $f$ in region $r$ is $\textit{evom}_{fr}$:
\begin{equation}
\textit{evom}_{fr} = \left\{
\begin{array}{lr}
\sum_j \textit{dfm}_{fjr} & f \in m_f \\
\textit{evom}_{fr} \textit{FT}_{fr} & f \in s_f
\end{array} \right.
\end{equation}
where the demand for primary factor is given by:
$$\textit{dfm}_{fjr}  = Y_{jr} \textit{vfm}_{fjr}
\left( \frac{c^f_{jr}}{p^{pf}_{fjr}}\right)^\sigma\,.$$

\paragraph{\bf Specific factors (\texttt{PS(f,j,r)})}

The net value of benchmark payments to factor $f$ in sector $j$ in region $r$ is \texttt{vfm(f,j,r)}:
\begin{equation}
\textit{vfm}_{fjr} \left( \frac{p^S_{s_{f},jr}}{p^F_{s_{f},r}} \right)^\eta =
\textit{dfm}_{fjr}\end{equation}
where the demand for primary factor is written above.

\subsubsection[Regional budget]{Regional budget (\texttt{RA(r)})}

Private and public incomes are given by :
\begin{equation}
\textit{RA}_r = \sum_f p^F_{fr} \textit{evom}_{fr} + p^I_n
\textit{vb}_r -  p^I_{r} \textit{vom}_{Ir} - p^G_{r}
\textit{vom}_{Gr}  + \mathcal{R}_r
\end{equation}
The base year current account deficit in region $r$ is \texttt{vb(r)}, and region $r=n$ corresponds to the ``numeraire region'' whose investment price (\texttt{P("i",rnum)} in the code) denominates international capital flows.
The numeraire region is the modelled region with the largest consumption (macro \texttt{loadrm}); in a single-region small open economy closure it is the modelled region itself, so that the current account is valued at its domestic investment price, the rest of the world providing the corresponding investment goods (following conventional static trade theory, we hold the current account deficit fixed in counterfactual analysis).
Furthermore, tax revenue in region $r$ consists of output taxes, intermediate demand taxes, factor taxes, final demand taxes, import tariffs, and export subsidies:
\begin{equation}
\begin{array}{ll}
\mathcal{R}_r = & \sum_j \mathcal{R}^o_{jr}
                 + \sum_{ij} \left( \mathcal{R}^{fd}_{ijr} +
                 \mathcal{R}^{fi}_{ijr} \right)
                 + \sum_{fj} \mathcal{R}^f_{fjr} \\
                &+ \sum_{i}  \left( \mathcal{R}^{pd}_{ir} +
                 \mathcal{R}^{pi}_{ir} + \mathcal{R}^{gd}_{ir} +
                 \mathcal{R}^{gi}_{ir}  \right)
                 - \sum_{is} \mathcal{R}^{xs}_{irs}
                 + \sum_{is} \mathcal{R}^{ms}_{isr}
\end{array}
\end{equation}
Each of these components of tax revenue can be calculated as an ad-valorem or proportional tax rate times a market price times the quantity demanded or produced.

Taxes related to $Y_{jr}$ include output taxes:\footnote{Tax revenues in the GAMS codes -- both \MCP and \MGE are represented by the \textit{macros} indicated in square brackets.}
$$\mathcal{R}^o_{jr} =  t^o_{jr} \textit{vom}_{jr} p^Y_{jr}
Y_{jr}\, \quad [\footnotesize{\texttt{REV\_TO(g,r)}}]\,,$$ tax revenue from intermediate inputs:
$$\mathcal{R}^{fd}_{ijr} =  t^{fd}_{ijr} p^Y_{ir}
\textit{ddfm}_{ijr}\, \quad [\footnotesize{\texttt{REV\_TFD(i,j,r)}}]\,,$$
$$\mathcal{R}^{fi}_{ijr} = t^{fi}_{ijr} p^M_{ir}
\textit{difm}_{ijr}, \quad [\footnotesize{\texttt{REV\_TFI(i,j,r)}}]\,,$$ and factor tax revenue:
$$\mathcal{R}^f_{fjr} = t^{f}_{fjr} p^F_{fr} \textit{dfm}_{fjr} \quad
[\footnotesize{\texttt{REV\_TF(f,g,r)}}]\,.$$

Taxes on household consumption of domestic and imported goods are:
$$\mathcal{R}^{pd}_{ir} = t^{pd}_{ir} p^Y_{ir} \textit{ddfm}_{iCr},
\quad [\footnotesize{\texttt{REV\_TFD(i,"C",r)}}]\,,$$ and
$$\mathcal{R}^{pi}_{ir} = t^{pi}_{ir} p^M_{ir} \textit{difm}_{iGr}
\quad [\footnotesize{\texttt{REV\_TFI(i,"C",r)}}]\,.$$

Taxes on public demand for domestic and imported goods are:
$$\mathcal{R}^{gd}_{ir} = t^{gd}_{ir} p^Y_{ir} \textit{ddfm}_{iGr}
\quad [\footnotesize{\texttt{REV\_TFD(i,"G",R)}}]\,,$$ and
$$\mathcal{R}^{gi}_{ir} =  t^{gi}_{ir} p^M_{ir} \textit{difm}_{iGr}
\quad [\footnotesize{\texttt{REV\_TFI(i,"G",r)}}]\,.$$

Export subsidies (paid by the government in the exporting region) are:
$$\mathcal{R}^{xs}_{irs} = t^{xs}_{irs} p^Y_{ir}
\textit{dxmd}_{irs} \quad [\footnotesize{\texttt{REV\_TXS(i,r,s)}}]\,,$$ and import tariff revenues are given by:
$$\mathcal{R}^{ms}_{isr} = t^{ms}_{isr} \left(p^Y_{is}
(1-t^{xs}_{isr}) \textit{dxmd}_{isr} + \sum_j p^T_j \textit{dtwr}_{jisr} \right) \quad [\footnotesize{\texttt{REV\_TMS(i,s,r)}}]\,.$$
